\section{НАЗВАНИЕ РАЗДЕЛА}

В процессе ручного контроля проектной документации, проверяющий выполняет достаточно рутинную работу,
для которой характерно возникновениe различных типовых ошибок.

\subsection{Подраздел с иллюстрацией}

На данной странице есть пример иллюстрации (рисунок~\ref{fig:example}).

\begin{figure}[htbp]
  \centering
  \includegraphics[width=0.85\textwidth]{test}
  \caption{Пример иллюстрации}
  \label{fig:example}
\end{figure}

\subsection{Название подраздела с таблицей и формулой}

В таблице~\ref{tab:example} приведены примеры использования.

\begin{table}[htbp]
  \caption{Пример таблицы}
  \label{tab:example}
  \begin{tabularx}{\textwidth}{|>{\centering\arraybackslash}p{0.18\textwidth}|>{\centering\arraybackslash}p{0.20\textwidth}|Y|}
    \hline
    Параметр & Номер использования & Использование \\
    \hline
    A & 1 & import документов правил \\
    \hline
    B & 2 & проверить проектный документ на соответвие правилам языка (пунктуация, орфография и тд) errors \\
    \hline
  \end{tabularx}
\end{table}

\Needspace{14\baselineskip}
\subsubsection{Пример таблицы без переноса}
Иногда могут появляться своеобразные таблицы с разными колонками.
\begin{table}[htbp]
  \caption{Пример построения таблицы}
  \label{tab:layout}
  \begin{tabularx}{\textwidth}{|Y|Y|Y|Y|}
    \hline
    \multirow{2}{*}{Заголовок} & \multicolumn{2}{c|}{Заголовок колонки} & \multirow{2}{*}{Заголовок} \\
    \cline{2-3}
    & Подзаголовок & Подзаголовок & \\
    \hline
    & & & \\
    \hline
  \end{tabularx}
\end{table}

\Needspace{18\baselineskip}
\subsubsection{Пример таблицы с переносом}

Достаточно часто используют перенос таблицы, как это случилось с таблицей 1.5.

\begin{tabularx}{\textwidth}{|>{\centering\arraybackslash}p{0.12\textwidth}|Y|Y|Y|}
\caption{Пример переноса таблицы}\label{tab:long}\\
\hline
Номер & Колонка & Колонки & Колонку \\
\hline
1 & 2 & 3 & 4 \\
\hline
\endfirsthead

\multicolumn{4}{r}{Продолжение таблицы~\thetable}\\
\hline
Номер & Колонка & Колонки & Колонку \\
\hline
1 & 2 & 3 & 4 \\
\hline
\endhead

\hline
\endfoot

\hline
\endlastfoot

A & данные & данные & данные \\
\hline
B & данные & данные & данные \\
\hline
C & данные & данные & данные \\
\hline
D & данные & данные & данные \\
\hline
E & данные & данные & данные \\
\hline
F & данные & данные & данные \\
\hline
G & данные & данные & данные \\
\hline
H & данные & данные & данные \\
\hline
I & данные & данные & данные \\
\hline
J & данные & данные & данные \\
\hline
J & данные & данные & данные \\
\hline
J & данные & данные & данные \\
\hline
J & данные & данные & данные \\
\hline
J & данные & данные & данные \\
\hline
J & данные & данные & данные \\
\hline
J & данные & данные & данные \\
\hline
J & данные & данные & данные \\
\hline
в & данные & данные & данные \\
\hline
J & данные & данные & данные \\
\hline
J & данные & данные & данные \\
\hline
J & данные & данные & данные \\
\hline
J & данные & данные & данные \\
\hline
J & данные & данные & данные \\
\hline
J & данные & данные & данные \\
\hline
\end{tabularx}

\subsubsection{Формулы}

Допускается нумерация в пределах главы
(раздела), в этом случае номер состоит из номера раздела и порядкового номера
формулы, разделенных точкой. 

\VKRformula[eq:example]{%
  E = \int_{min}^{max}\log_{234234}^{12312}\frac{mc^2}{mc^4},
}{%
  \VKRwitem{$E$}{значение энергии}%
  \VKRwitem{$m$}{масса тела}%
  \VKRwitem{$c$}{скорость света}[.]%
}

Пример расшифровки символов представлен в единственной формуле (1.1) в этом шаблоне.

